\section{Consequence for the Cascade Exponent $\beta^*$}
  \label{sec:beta}

  \subsection{The structural relation $\delta = f(\beta)$: three levels}

    The capacity exponent $\delta$ and the cascade exponent $\beta^*$ are linked by a
    structural relation derived in the O3--O7
    sub-programme~\cite{Beau2026a7,Beau2026a8,Beau2026a9,Beau2026a10,Beau2026a11}.
    We use this relation at three distinct levels, which must not be conflated.

    \textbf{Level 1 (empirical).}
    The present paper provides the measurement: $\hat{\delta}_{\mathrm{exact}}\approx X$
    with uncertainty $\pm\sigma$ from Table~\ref{tab:delta}.

    \textbf{Level 2 (structural).}
    The O3--O7 papers establish the mapping $\delta\mapsto\beta^*$ from first principles
    within the Cosmochrony framework, without fitting any free parameter to the lepton data.
    We do not refit this relation here.
    Rather, we propagate $\hat{\delta}_{\mathrm{exact}}$ through the pre-existing structural
    relation as a change of input.

    \textbf{Level 3 (phenomenological).}
    O3~\cite{Beau2026a7} established the window $\beta^*\in(0.09,0.13)$ as the range
    compatible with the observed charged lepton mass ratios
    $m_e:m_\mu:m_\tau\approx 1:207:3477$.
    This window was derived independently of $\delta$; it serves here as a consistency
    target.

  \subsection{Propagating $\hat{\delta}_{\mathrm{exact}}$ through the O3--O7 relation}

    The structural relation between the capacity exponent $\delta$ and the cascade
    exponent $\beta^*$ is derived in O7~\cite{Beau2026a11} from the subdiffusive
    valence growth law established in O4~\cite{Beau2026a8}.
    In the notation of those papers, $\delta$ plays the role of the capacity decay
    exponent $\alpha$ in $\Sigma_n \sim p(n)^{-\alpha}$, and the relation is
    \begin{equation}
      \beta^* = g(\delta) = \frac{1}{\delta + \tfrac{1}{2}},
      \qquad\text{equivalently}\qquad
      \delta = f(\beta^*) = \frac{1}{\beta^*} - \frac{1}{2}.
      \label{eq:delta-beta-relation}
    \end{equation}
    This relation is not fitted to the lepton data in the present paper; it is used
    here as a fixed structural output of O3--O7, with $\hat{\delta}_{\mathrm{exact}}$
    as the only new input.
    The O7 paper also records the consistency check: the window
    $\alpha \in [7.4, 10.6]$ (the range anticipated from O10 at the proxy level)
    corresponds precisely to $\beta^* \in (0.09, 0.13)$ via~\eqref{eq:delta-beta-relation},
    confirming internal consistency of the O3--O7 chain.

    Substituting $\delta = \hat{\delta}_{\mathrm{exact}} \pm \sigma$ into~\eqref{eq:delta-beta-relation}
    yields the exact-projection constraint on $\beta^*$ reported in Table~\ref{tab:beta}.

    \begin{table}[ht]
      \centering
      \caption{Constraint on $\beta^*$ derived from the exact-projection measurement
      via the structural relation~\eqref{eq:delta-beta-relation}.
      The phenomenological target window $\beta^*\in(0.09,0.13)$ is shown for comparison.
      All $\beta^*$ values lie significantly above the target window: incompatible outcome.}
      \label{tab:beta}
      \begin{tabular}{ccccc}
        \hline
        $q$ & blocs & $\hat{\delta}_{\mathrm{exact}}$ & $\beta^* = 1/(\hat{\delta}+\tfrac{1}{2})$ &
        Compatible with $(0.09,0.13)$? \\
        \hline
        29                    & 50 $\times$ 3\,seeds & $4.83 \pm 0.24$ & $0.184\text{--}0.193$ & No            \\
        53                    & 20                   & $4.40$          & $0.204$               & No            \\
        61                    & 15                   & $4.57$          & $0.197$               & No            \\
        \hline
        O11 proxy ($q{=}307$) & 500                  & $3.39$          & $0.257$               & No            \\
        O7 anticipated        & ---                  & $[7.4,10.6]$    & $[0.09,0.13]$         & \emph{Target} \\
        \hline
      \end{tabular}
    \end{table}

  \subsection{Consistency check and interpretation of tension}

    Table~\ref{tab:beta} shows a uniform \textbf{incompatible outcome}: all primes yield
    $\beta^* \approx 0.19$--$0.20$, lying significantly above the phenomenological target
    $\beta^* \in (0.09, 0.13)$ from O3~\cite{Beau2026a7}.
    The gap is approximately a factor of 1.5--2 in $\beta^*$, corresponding to a gap
    of approximately 4--5 units in $\delta$ relative to the anticipated range
    $\alpha \in [7.4, 10.6]$ from O7.
    We report this incompatibility without adjustment.

    There are three structurally distinct sources of tension.

    \textbf{Source 1: insufficient prime range.}
    The primes tested ($q \leq 61$) are far below the asymptotic regime
    identified in O11 ($q \approx 211$--$307$).
    The exact-block fitting window is 4--7 data points at $q \leq 61$,
    compared to 22 points in O11 at $q = 307$.
    It is plausible that $\hat{\delta}_{\mathrm{exact}}$ continues to grow as $q \to \infty$,
    reaching the range $[7.4, 10.6]$ at $q \geq 211$.
    The $q$-dependence observed in the proxy series (where $\hat{\delta}_{\mathrm{cap}}$
    decreased from 5.78 at $q=101$ to 3.39 at $q=307$) shows that large prime corrections
    can be substantial; the exact series may exhibit a different but comparably large drift.
    This interpretation makes O12-O1 (extension to $q \geq 101$) the most important
    open computation in the programme.

    \textbf{Source 2: structural relation requiring revision.}
    The relation $\beta^* = 1/(\delta + \frac{1}{2})$ was derived in O3--O7 using the
    proxy-level dynamics as input at certain steps.
    If the exact-block dynamics differ structurally from the proxy dynamics
    (as the present paper establishes they do, via the $\gamma$-induced acceleration),
    then the derivation of $f(\beta)$ may need to be revisited with the exact observable
    as input.
    This would leave~\eqref{eq:delta-beta-relation} as an approximation valid at the
    proxy level, requiring a corrected relation $f_{\mathrm{exact}}(\beta)$ derived
    from the metaplectic structure of $\rho_c$.

    \textbf{Source 3: revision of the phenomenological window.}
    A third possibility is that the predicted window $\beta^* \in (0.09, 0.13)$ itself
    requires revision in light of additional physical inputs.
    This is the least preferred interpretation in the absence of independent motivation.

    The three sources are not mutually exclusive.
    The programme currently cannot discriminate between them without large-$q$ data.
    The incompatible outcome is therefore honest but not conclusive: it identifies a
    quantitative tension in the Cosmochrony programme that the next computational step
    (O12-O1, $q \geq 101$) is designed to resolve.
